\begin{figure}[t]
\centering
\begin{adjustbox}{max width=\textwidth,center}
\begin{tikzpicture}[
  node distance=7mm and 8mm,
  sector/.style={draw,rounded corners,fill=gray!7,align=center,
    minimum height=13mm,text width=0.25\textwidth,font=\small},
  map/.style={draw,very thick,rounded corners,align=center,
    minimum height=14mm,text width=0.22\textwidth,font=\small},
  result/.style={draw,double,rounded corners,align=center,
    minimum height=14mm,text width=0.22\textwidth,font=\small},
  arrow/.style={-{Latex[length=2.2mm]},thick}
]
\node[sector] (node) {nodes\\[-1mm]
  $t_a=h_a\widetilde h_a$\\
  wrapped-M2 hypers};
\node[sector,right=of node] (e2) {$E_2$ sectors\\[-1mm]
  $\beta_r=\mu^{\CC}_{SU(2)}$\\
  reduced Higgs line};
\node[sector,right=of e2] (e3) {$E_3$ sectors\\[-1mm]
  $(H_{p1},H_{p2})$ or $L_p$\\
  branch profile $\sigma$};

\node[map,below=12mm of e2] (charge) {compact charge map $Q_{X,\sigma}$\\[1mm]
  $c_{Ia}=D_I\!\cdot C_a$\,,\quad
  $m_{Ip\alpha}=D_I\!\cdot\iota_{p*}R_{p\alpha}$};
\node[result,right=15mm of charge] (kernel) {allowed holomorphic data\\[1mm]
  $\ker Q_{X,\sigma}$};

\draw[arrow] (node.south) -- ++(0,-5mm) -| ([xshift=-6mm]charge.north);
\draw[arrow] (e2.south) -- (charge.north);
\draw[arrow] (e3.south) -- ++(0,-5mm) -| ([xshift=6mm]charge.north);
\draw[arrow] (charge.east) -- node[above,font=\scriptsize]
  {$\mu_I^{\CC}=0$} (kernel.west);

\node[align=center,font=\scriptsize,below=5mm of charge,text width=0.66\textwidth]
  {The same divisor--curve pairing is the wrapped-M2 electric charge,
   the flavor-cocharacter gauging, and the dual of the geometric
   local-to-global map.};
\end{tikzpicture}
\end{adjustbox}
\caption{The compact Higgs-gluing mechanism for the sectors in the toric
examples, on a fixed $E_3$ branch profile.
The displayed kernel is the image in invariant complex moment-map
coordinates. Balanced representatives lift it to the zero level of the full
triplet moment maps. The exact geometric smoothing criterion requires this
kernel to meet the complement of every local discriminant.}
\label{fig:local-to-global}
\end{figure}
